\documentclass[10pt]{article}
\usepackage[T1]{fontenc}
\usepackage[utf8]{inputenc}
\usepackage[french]{babel}
\usepackage{amsmath,amssymb,mathrsfs,mathpazo}
\usepackage{geometry}
\geometry{a4paper,margin=2cm}
\usepackage{graphicx,array,tabularx,tabularray,colortbl,multirow,multicol,booktabs,enumitem,pifont,ulem,fancybox,diagbox,cancel}
\usepackage{pstricks,pstricks-add,pst-plot,pst-tree,pst-node,pst-text,pst-func,pst-math,pst-bspline,pst-3dplot,pst-eucl,pst-solides3d}
\let\pscalculate\undefined
\usepackage{pst-bezier}
\usepackage{tikz,pgfplots,tkz-tab}
\usetikzlibrary{arrows,arrows.meta,patterns,calc,positioning,trees,decorations.markings,shapes.geometric,shapes.misc,shapes.arrows,intersections,decorations.pathmorphing,decorations.pathreplacing,backgrounds,shadows,babel}
\pgfplotsset{compat=1.18}
\usepgfplotslibrary{fillbetween}
\usepackage[np]{numprint}
\usepackage{listings,makecell,pas-tableur,twoopt}
\let\barre\undefined
\usepackage[pstricks]{tablvar}
\definecolor{darkgreen}{RGB}{0,100,0}
\providecommand{\red}{\color{red}}
\providecommand{\blue}{\color{blue}}
\providecommand{\green}{\color{green}}
\providecommand{\black}{\color{black}}
\providecommand{\white}{\color{white}}
\providecommand{\gray}{\color{gray}}
\providecommand{\cyan}{\color{cyan}}
\providecommand{\magenta}{\color{magenta}}
\providecommand{\yellow}{\color{yellow}}
\providecommand{\violet}{\color{violet}}
\providecommand{\orange}{\color{orange}}
\providecommand{\euro}{\texteuro}
\providecommand{\EUR}{\texteuro}
\providecommand{\ieme}{\textsuperscript{e}}
\providecommand{\degres}{\ensuremath{^\circ}}
\providecommand{\degree}{\ensuremath{^\circ}}
\providecommand{\vv}[1]{\overrightarrow{#1}}
\providecommand{\Coord}[1]{#1}
\providecommand{\decosix}{}
\providecommand{\starredbullet}{\ensuremath{\bullet}}
\pagestyle{empty}
\setlength{\parindent}{0pt}
\usepackage{adjustbox}\DeclareUnicodeCharacter{2236}{:}\DeclareUnicodeCharacter{2190}{\ensuremath{\leftarrow}}
\providecommand{\R}{\mathbb{R}}
\providecommand{\N}{\mathbb{N}}
\providecommand{\Z}{\mathbb{Z}}
\providecommand{\Q}{\mathbb{Q}}
\providecommand{\D}{\mathbb{D}}
\providecommand{\C}{\mathbb{C}}
\providecommand{\e}{{\mathrm{e}}}
\providecommand{\vect}[1]{\overrightarrow{#1}}
\providecommand{\vectt}[1]{\overrightarrow{\mathrm{#1}}}
\providecommand{\vvt}[1]{\overrightarrow{\mathrm{#1}}}
\providecommand{\np}[2][]{{#2~#1}}
\providecommand{\vv}[1]{{\overrightarrow{#1}}}
\renewcommand{\d}{{\mathrm{d}}}
\providecommand{\degres}{{^\circ}}
\providecommand{\degre}{{^\circ}}
\providecommand{\blue}{\color{blue}}
\providecommand{\red}{\color{red}}
\providecommand{\enskip}{\quad}
\providecommand{\llbracket}{\lbrack\!\lbrack}
\providecommand{\rrbracket}{\rbrack\!\rbrack}
\providecommand{\euro}{{\text{€}}}
\providecommand{\no}{{\text{n°}}}
\providecommand{\interval}[2]{{[#1;#2]}}
\providecommand{\intervFO}[2]{{[#1;#2[}}
\providecommand{\dotfill}{\ldots}

\providecommand{\ds}{\displaystyle}
\providecommand{\textding}[1]{#1}
\providecommand{\Oij}{\(\mathrm{O};\vec{\imath},\vec{\jmath}\)}
\providecommand{\Oijk}{\(\mathrm{O};\vec{\imath},\vec{\jmath},\vec{k}\)}
\providecommand{\Ouv}{\(\mathrm{O};\vec{u},\vec{v}\)}
\providecommand{\intervFF}[2]{\ensuremath{[#1;#2]}}
\renewcommand{\euro}{\ifmmode\text{\texteuro}\else\texteuro\fi}
\DeclareUnicodeCharacter{00B0}{\ifmmode^{\circ}\else\textdegree\fi}
\DeclareUnicodeCharacter{20AC}{\ifmmode\text{\texteuro}\else\texteuro\fi}
% Énoncés : Éducation nationale. Transcription et corrigés : APMEP, Denis Vergès et contributeurs. Découpage et classement : CoopMaths, Sébastien Lozano, Éric Elter et contributeurs. Adaptation de la présentation : Maths Explorer.
% CC BY-SA : https://coopmaths.fr/www/about/
% https://coopmaths.fr/alea/static/e3c/2020/tex/e3c_2020_00_sujet21_4_cor.tex
\begin{document}
\noindent\textit{APMEP / CoopMaths. CC BY-SA. Présentation adaptée par Maths Explorer.}\par\medskip
\subsection*{1.}
	
	Dans le service \(A\), il y a 450 personnes sur un effectif de 1 000, soit une proportion de
	\[
	p(A) = \dfrac{450}{1000} = 0,45.
	\]
	
	Parmi ces personnes du service \(A\), 40 \% résident à moins de 30 minutes de l’entreprise, donc
	\[
	p_A(T) = \dfrac{40}{100} = 0,4.
	\]
	
	\subsection*{2.}
	
	\begin{center}
	\begin{tikzpicture}
		[level 1/.style={level distance=3cm,
			sibling distance=2.2cm},
		level 2/.style={level distance=3cm,
			sibling distance=1.5cm},
		]
		\node {} [grow'=right]
		child {node {$A$}
			child {node {$T$}
				edge from parent node[above] {$0,4$}
			}
			child {node {$\overline T$}
				edge from parent node[below] {$0,6$}
			}
			edge from parent node[left] {$0,45$}}
		child {node {$B$}
			child {node {$T$}
				edge from parent node[above] {$0,2$}
			}
			child {node {$\overline T$}
				edge from parent node[below] {$0,8$}
			}
			edge from parent node[above] {$0,23$}}
		child {node {$C$}
			child {node {$T$}
				edge from parent node[above] {$0,8$}
			}
			child {node {$\overline T$}
				edge from parent node[below] {$0,2$}
			}	edge from parent node[left] {$0,32$}}
		
		;
	\end{tikzpicture}
\end{center}
	
	\subsection*{3.}
	
	Il faut trouver \(p(A \cap T) = p(A) \times p_A(T) = 0,45 \times 0,4 = 0,18\).
	
	\subsection*{4.}
	
	On a de même \(p(B \cap T) = p(B) \times p_B(T) = 0,23 \times 0,2 = 0,046\), puis
	\[
	p(C \cap T) = p(C) \times p_C(T) = 0,32 \times 0,8 = 0,256.
	\]
	
	D’après la loi des probabilités totales :
	\[
	p(T) = p(A \cap T) + p(B \cap T) + p(C \cap T) = 0,18 + 0,046 + 0,256 = 0,482.
	\]
	
	\subsection*{5.}
	
	Il faut trouver \(p_T(C) = \dfrac{p(T \cap C)}{p(T)} = \dfrac{p(C \cap T)}{p(T)} = \dfrac{0,256}{0,482} \approx 0,5311\), soit 0,531 au millième près.

\end{document}
